\section{Conclusion}
\label{sec:conclusion}
This review set out to examine surrogate-assisted energy system
optimization with a particular focus on multi-energy systems and
multi-objective settings. The main result is not a single best model
class, but a clearer picture of how the field is structured. Model class
alone is not enough to explain the literature. The same Gaussian
process, neural network or polynomial surrogate can serve different
purposes depending on what it approximates, where it is inserted into
the optimization workflow, and how its output is validated. For that
reason, the review organized the field along four linked dimensions:
model class, surrogate target, integration depth and trust mechanism.
This was complemented by a synthesis of training-data generation,
design-of-experiments choices, integration patterns, validation
practice, and application evidence.
Three observations stand out across the reviewed studies. First,
surrogates are already used in several distinct ways rather than in one
generic pattern. Some studies replace expensive inner evaluations,
others accelerate search, provide warm starts, support decomposition, or
handle uncertainty. The dominant pattern depends on the underlying
decision problem. Second, substantial computational savings are reported
across many applications, but the supporting evidence is uneven. Runtime
reductions are usually shown more clearly than their effect on decision
quality. Third, the most active application area is multi-objective
design in microgrids, hybrid renewable systems and broader MES
configurations, while stronger validation practice is more common in
dispatch, unit commitment and OPF studies.
The evidence map also shows where the field is still thin. Publicly
reusable datasets are scarce, especially for district heating and
sector-coupled MES cases. Validation is often still centered on
predictive fit rather than on feasibility, regret or ranking stability.
Many studies mix surrogate modeling, optimization and domain
engineering effectively at case-study level, but the reporting style
does not yet make those studies easy to compare or reuse. This is where
the six challenge areas discussed above become important: common
benchmarks, calibrated uncertainty, better handling of distribution
shift, reproducible datasets, explicit treatment of uncertainty in MCDM,
and less fragmented toolchains.
Taken together, the review provides a structured and quantitative
overview of surrogate-assisted energy system optimization. It contributes
a taxonomy for organizing the literature, an evidence map for tracing
connections between methods and applications, and a more concrete basis
for discussing where surrogate models already work well and where the
current evidence is still weak. If the field moves toward clearer
benchmarks, better reporting and more reusable case-study artefacts,
surrogate-assisted optimization can become easier to compare, easier to
reproduce and more useful for practical MES decision support.
